\section{Asymptotics of the two color sums: the angular calculation}\label{c4:angular}

The remaining one-mark estimate concerns the ratio of the decorated and undecorated color sums.  We first express both as angular traces; the radial transformation will then isolate the different powers contributed by their long central intervals.


Here and for the next (radial and spectral) calculations we change notation. Matrices, charges and scalar products below are auxiliary algebra; e.g. \(C,A,h,v\) below do not denote a cup/fusion tensor or vectors in physical space. The quantities to which we will apply the calculation are the unmarked \(N\)-site color sum \(\tau_N\) and the marked sum \(\tau^v_{N-1}\) (with \(N-1\) color sites and the one distinguished parameter, no factor \(1/2\)), at homogeneous parameters. Recall \(q=e^{\pi i/4}\), and write the three color weights here as
\[
 w_-=\frac{-1-\sqrt{2}+i}{2},\qquad w_0=2+\sqrt{2},\qquad w_+=q w_- .
\]
Homogeneous values of the color sums may be taken by analytic confluence as in the color identities. We derive an expression for them via elementary oscillator traces, then transform the expression. A convergence point (real-analytic continuation of the transformed series) used at the end of the angular calculation is proved separately below by the radial estimates.

\subsection{A formal charge space and root operators}

Let \(\alpha_j\), \(j\in\mathbb Z/8\), be the images of the coordinate vectors in \(\mathcal V=\mathbb Q^8/\mathbb Q(1,\ldots,1)\), and let \(\mathcal L\) be their integral span. Use base extensions to \(\mathbb R,\mathbb C\) when needed. A symmetric bilinear form \((\, ,\,)\) on this quotient has circulant Gram matrix with first row
\[
 (2,-1,1,-2,2,-2,1,-1).
\]
It is nonsingular (the principal size-7 minor is 24) but not positive definite. Write \(b^2=(b,b)\), \(M^*\) for bilinear adjoints, and \(\mathcal L^*=\{b\in\mathcal V:(b,\mathcal L)\subset\mathbb Z\}\). Put \(C\alpha_j=\alpha_{j+1}\), \(K=(C-I)^{-1}=\sum_{j=0}^7 j C^j/8\), so \(K^*=-I-K\). Define \(h,v\) by their pairings in order with \(\alpha_j\):
\[
 h:(1,0,1,0,0,-2,0,0),\qquad v:(0,1,1,1,-1,-1,-1,0),\qquad p=v-h.
\]
Let \(r=(\alpha_0,0,-\alpha_7,-\alpha_1)\), \(e_{ij}=r_i-r_j\) (labels \(1,2,3,4\)), \(h_i=h+e_{i1}\). Put
\[
 d_0=\alpha_6+\alpha_1,\quad d_1=Cd_0,\quad \beta=e_{34},\qquad
 S_b x=x-(x,b)b\ \ (b^2=2),\quad A=S_{\alpha_1} C .
\]
We record linear computations for use here and below.
\par\noindent\textbullet\quad \((\alpha_0,K\alpha_j)=(-1,1,0,1,-1,1,-1,0)_j=(p,\alpha_j)\), \(K\alpha_1=-p\), and \(K(\mathcal L+\mathbb Z h)\subset\mathcal L^*\).
\par\noindent\textbullet\quad The \(e\)'s have the usual pairings of coordinate differences in Euclidean four-space; \((e_{ij},h_k)=\delta_{ik}-\delta_{jk}\). Both \(d_j\) have square 0 and are orthogonal to every \(e_{ab}\); \((d_0,h)=(d_0,v)=0\). Also \(h^2=h_i^2=8/3\), \(p^2=9/8\), \((p,h)=0\), and \((v,e_{i4})=1\) for \(i<4\).
\par\noindent\textbullet\quad \(Ap=p,\ (A+A^{-1})h=h,\ \det(X-A)=(X-1)(X^2-X+1)(X^2+X+1)^2\).
Indeed one can multiply the circulant and shift above using representatives \(h=(13,5,13,5,-11,-19,-11,5)/24,\ v=(5,31,41,19,-19,-41,-31,-5)/48\); \((\alpha_j,Kh)=(0,0,-1,-1,-1,1,1,1)_j\). For the determinant use \(A=C-\alpha_1(\alpha_0,\cdot)\), \((X-C)^{-1}=(X^8-1)^{-1}\sum_{j=0}^7 X^{7-j} C^j\) in the rank-one formula.

The algebraic construction belongs to the lattice vertex-operator method of Frenkel--Kac and Segal~\cite{FrenkelKac1980,Segal1981}; see also \cite[Section~2.1]{Frenkel1985}. We specify the charge lattice and cocycle and derive the root commutators by normal ordering. The ambient bilinear form here is indefinite, so the initial fields and traces are interpreted coefficientwise. The finite-support argument below justifies their algebraic rearrangement; analytic convergence is established separately.

Take as states a copy for each \(n\in\mathcal L^*\) of the polynomial symmetric algebra on copies of \(\mathcal V_{\mathbb C}\) indexed by \(k=1,2,\ldots\) (algebraic direct sum of the momentum summands). Operators \(P,T_b\) are respectively vector momentum (value \(n\)) and shift by \(b\). Write \(c_{-k}(x)\) for multiplication by \(x\) in copy \(k>0\), \(c_k(y)\) for its contracted derivative, with \([c_k(y),c_{-k}(x)]=k(y,x)\); distinct copies commute. Define \(H=P^2/2+\) the weighted polynomial degree (copy \(k\) has weight \(k\)). For \(b\in\mathcal L^*,\ b-\sigma\in\mathcal L+\mathbb Z h\) (we use \(\sigma=0,\pm p\)), use fields
\[
 W_b^\sigma(z)=T_b z^{(b,P)+b^2/2} e^{\pi i(b,KP)+2\pi i(P,K\sigma)}
       \exp\!\left(\sum_{k>0} c_{-k}(b)z^k/k\right)
       \exp\!\left(-\sum_{k>0} c_k(b)z^{-k}/k\right).
\]
Powers have bounded denominators and output is read power by power (exponent equals change in \(H\)); such a coefficient is an algebraic operator, also row-finite on a monomial basis. For \(a\in\mathcal L,\ a^2=2\), put \(X_a=[z^0]W_a^0(z)\); this preserves \(H\), and is locally nilpotent since it shifts by \(a\) at fixed \(H\) (also locally nilpotent backwards). The commutators needed are
\[
 \begin{aligned}
 [X_a,W_b^\sigma(z)]&=
 \begin{cases}0 &(a,b)\ge0,\\ e^{\pi i(a,K b)} W_{a+b}^\sigma(z)&(a,b)=-1,\end{cases}\\
 [X_a,W_{-a}^0(z)]&= e^{-\pi i(a,K a)}\left[(a,P)+\sum_{k\ne0}c_k(a)z^{-k}\right].
 \end{aligned}
\]
Indeed order \(W_a^0(\zeta) W_b^\sigma(z)\) with shifts at the far left and creators before annihilators; the extra scalar from the shift and contractions is \(e^{\pi i(a,K b)}(\zeta-z)^{(a,b)}\) expanded for \(\zeta\) outside. Reversing gives the same rational power expanded inside, with matching phase because \((a,K(b-\sigma))\in\mathbb Z\). Thus the zero-mode commutator is the residue at \(\zeta=z\) with measure \(d\zeta/\zeta\), giving the formulas (one derivative for the double pole). This residue check is valid between specified input and output monomials where the remaining normal factor is Laurent in \(\zeta\). Also \([c_{-k}(y),W_a^0(z)]=(a,y)z^{-k}W_a^0(z)\).

Normalize, for \(i\ne j\),
\[
 E_{ij}=e^{-\pi i(e_{ij},K h_j)} X_{e_{ij}} .
\]
Together with the diagonal differences \((e_{ij},P)\) these represent \(\mathfrak{sl}_4\); they act on the fields \(W_{h_l}^0\) by commutator exactly as matrix units and Cartan differences on column vectors. Indeed the brackets of nonopposite root operators with pairing \(-1\) are scalar multiples of the sum-root operator by the formula, with scalar determined on the \(W_{h_l}^0\); opposites give a multiple of \((e_{ij},P)\) the same way. The action integrates to \(SL_4(\mathbb C)\): each vector is in a finite-dimensional stable subspace by fixed \(H\) and positivity on the \(e_{ij}\) span. Recall briefly that for a finite-dimensional bracket representation one can integrate \(g^{-1}dg\) in the representation along paths (transport \(dU=U\,d\rho(g^{-1}dg)\)); differentiating along a homotopy gives endpoint independence by the bracket identity. The group is simply connected: retract to \(SU_4\) by polar decomposition, contract a loop of first columns on the unit sphere and lift this homotopy starting at the given frames, using local orthonormal sections with determinant adjustment (uniform small steps along the homotopy and last frames as origin). Then contract in \(SU_3\) and repeat. Equivalently this works piecewise smoothly (a smooth curve on a sphere of dimension \(>1\) misses a point). We use hats for group operators or just their matrices when unambiguous. Let \(\widehat C\) send momenta and oscillator factors by \(C\).

\subsection{Rotating the three-color insertion}

We now choose a finite root rotation that replaces the three-color insertion by one charge. The price is a finite list of integrated root fields; their charges will determine the coercive radial interaction.

Set
\[
 w=-e^{\pi i/8},\quad c=w^{-2},\qquad
 G_0=\begin{pmatrix}1&0&0&-c\\ w&1&0&0\\ w^2&w&1&0\\0&0&0&1\end{pmatrix},
 \qquad g=\widehat G_0\exp(t X_{\beta-d_0}),
\]
where \(t\) will be chosen shortly. Conjugating \(W_h^0\) by \(g\) gives
\[
 \widetilde T=W_h^0+w W_{h_2}^0+w^2 W_{h_3}^0.
\]
Conjugating \(W_v^p\) leaves it unchanged. In fact \((\beta-d_0)^2=2,\ (\beta-d_0,h)=0,\ (\beta-d_0,v)=1\), and \(G_0\) uses the \(\mathfrak{sl}_3\) on \(123\) and the \(i4\) columns (\(i<4\)). Conjugating \(W_{-v}^{-p}\) by \(g^{-1}\) gives a fixed finite linear combination, called the rotated compensator below, of fields \(W_{b_*}^{-p}\), \(b_*\in-v+\mathcal L,\ b_*^2=v^2\). Indeed the finite part acts within the four fields of weights \(-v+e_{i4}\); the remaining root has pairing \(\ge-1\) and string length at most one more step from these fields. These statements concern operators, power by power, not just projected matrix entries.

We have the exact factorization
\[
 \begin{split}
 g^{-1}\widehat C g &= \widehat U\, D\, \widehat S\,\widehat C,\\
 D&=\exp((-7\pi i/8)(e_{42},P)),\qquad
 \widehat S=\exp E_{42}\exp(-E_{24})\exp E_{42},
 \end{split}
\]
where \(\widehat U\) is a fixed finite product of factors \(\exp(u' X_b)\), \(u'\) complex constants, with \(b\)'s from
\begin{equation}
 \beta-d_0,\quad e_{12},\quad e_{14},\quad e_{42},\quad
 e_{32}+d_1,\quad \beta,\quad \beta+d_1 .                                      \label{c4:Ang-1}
\end{equation}
Here is a check, to avoid assuming any general factorization for these operators. The first factor of \(\widehat U\) is \(g^{-1}\widehat G_0=\exp(-t X_{\beta-d_0})\). For the rest use matrix notation, with the row 3 columns \(j\ne3\) also allowed terms linear in \(z_1\), a symbol for shift by \(d_1\). More precisely, the invariant finite subgroup fixing column 3 (with determinant 1 on \(124\)) normalizes the extra commuting shifted row copy as it does an ordinary row 3. Let \((\pi(1),\pi(2),\pi(3))=(4,2,1)\) and \(l=(1,-1,1)\). Conjugation by \(\widehat C\) sends \(\mathfrak{sl}_3\) on \(123\) to negative transpose on the \(\pi\)-labels conjugated by \(\operatorname{diag}(l)\); this follows from \(C e_{ij}=e_{\pi(j),\pi(i)}\) there and the phase table
\[
 \big[-(e_{ij},K h_j)\big]_{ij}
 =\begin{pmatrix}0&-1&-2&-1\\0&0&-1&-1\\1&0&0&-1\\0&0&0&0\end{pmatrix}.
\]
Since \(C e_{i4}=e_{3,\pi(i)}+d_1\) for \(i<4\), normalize the shifted row by declaring \(\widehat C E_{i4}\widehat C^{-1}\) to represent \(-l_i E_{3,\pi(i)} z_1\). Its bracket action with \(\mathfrak{sl}_3\) on \(124\) is the indicated row action by conjugation, and shifted-row elements commute with finite row 3 by nonnegative pairings, and with each other. Thus this realizes the semidirect row matrix identities exactly, exponentiating commuting locally nilpotent rows transforming by the finite group. Finally \(C(\beta-d_0)=e_{31}\). Using the inverse-transpose and row conversion on \(G_0\), the matrix for \(\widehat C g\widehat C^{-1}\) is
\[
 \begin{pmatrix}
 1&0&0&0\\ w&1&0&0\\ a&c w z_1&1&c z_1\\0&w&0&1
 \end{pmatrix}
\]
with arbitrary \(a\) (vary \(t\) multiplying by an ordinary \(31\) factor on the right). Choose \(a=w^2\). The remaining factor \(\widehat G_0^{-1}(\widehat Cg\widehat C^{-1})\widehat S^{-1}D^{-1}\) is then
\[
 \begin{pmatrix}1&-1/w&0&c\\0&1&0&0\\0&-z_1/w&1&c z_1-1\\0&-w&0&1\end{pmatrix},
\]
which is a product using \eqref{c4:Ang-1}: first the \(42\) transvection on the left, then the entries in rows \(1,3\) columns \(2,4\).

The operator \(\widehat S\) acts on a pure momentum state \(n\) by \(S_{\alpha_1}\) with scalar \(e^{-\pi i k(\alpha_1,K n)}\), \(k=(\alpha_1,n)\). It sends creation operators by conjugation according to the same reflection without the scalar. For the latter, \(c_{-j}(\alpha_1)\) is in an adjoint \(\mathfrak{sl}_2\) triple with the corresponding field modes of the two roots by the commutators, and orthogonal creators commute. For the scalar, a pure momentum state with \(k\ge0\) is highest for \(E_{24}\) by degree, and \(\widehat S\) sends it to \(E_{42}^k/k!\) on it by the usual raising/lowering string (or the degree-\(k\) symmetric power of the two-dimensional representation). At the output charge \(S_{\alpha_1}n\) there can be no nonempty oscillators at the given \(H\). Here \(E_{42}=X_{-\alpha_1}\); multiplying its \(k\) fields in order between empty oscillators produces cocycle \(e^{-\pi i k(\alpha_1,K n)}(-1)^{\binom{k}{2}}\) times
\(\prod z_j^{1-k}\prod_{j<l}(z_j-z_l)^2\), whose constant term is \((-1)^{\binom{k}{2}}k!\) by the two Vandermondes. For \(k<0\) use the lowest state and \((-E_{24})^{-k}=X_{\alpha_1}^{-k}\) (power here of positive integer degree) in the same way. Hence \(D\widehat S\widehat C\) acts by \(A\) on oscillator factors, and sends momentum \(m\) to \(n=A m\) with multiplier
\[
 \exp(2\pi i(\mu,n)-\pi i(\alpha_0,m)(\alpha_0,K m)),\qquad \mu=-7e_{42}/16 .
\]

\subsection{Angular traces and their finite support}

Fix real \(L>0\) and now put \(Q_L=e^{-4\pi^2/L}\), \(s=u+i\phi\), \(z(s)=e^{-2\pi i s/L}\) (powers use this log). Use vertices \(W_b^\sigma(z(s))\) scaled by \((2\pi/L)^{b^2/2}\); ``earlier to later'' angularly means increasing \(\phi\) in an interval of width \(<2\pi\). Also put \(J=e^{-\pi i P^2}\), commuting with all the zero-mode rotations used above by lattice integrality. Consider the angular trace
\[
 \operatorname{Tr}\!\left[J\widehat C Q_L^H (\text{scaled vertices, latest first})\right]/Z_C,\qquad
 Z_M=\prod_{j\ge1}\det(1-M Q_L^j)^{-1}.
\]
In the unmarked case \(\epsilon=0\) take \(N\) copies of scaled \(\widetilde T\). In the marked case \(\epsilon=1\) take chronologically scaled \(W_{-v}^{-p}\) at \(s_*=u_*+i\phi_*\), with \(0<u_*<L,\ \phi_*=-\pi/3\), then \(W_v^p\) and \(N-1\) copies of \(\widetilde T\) (scaled). The cluster of \(N\) vertices is always at \(s_j=x\lambda_j\), \(x\to0^+\), with fixed \(\lambda_j\) of strictly increasing imaginary parts and decreasing real parts. Divide the trace by \(\prod_{j<k\ {\rm in\ cluster}}(s_j-s_k)^{h^2}\) (principal logs) and take this limit, calling the result the confluent trace.

Here is its transformed angular expression, equal up to a constant unit factor (depending on fusion conventions and \(N,\epsilon\)). Replace \(\widehat C\) in the twist by \(D\widehat S\widehat C\) (not in \(Z_C\)), and take as scaled vertices, now without the cluster division:
\par\noindent\textbullet\quad factors from \(\widehat U\) in rightmost-first order, expanding each \(\exp(u' X_b)\) as a sum over numbers \(k_b\ge0\) of copies of the scaled vertex \(W_b^0\) integrated over \(u\in[0,L]\), measure per copy \(u'\,du/(2\pi)\) and divisor \(k_b!\). Each factor has its own contour angle \(\phi\); use distinct angles close to \(-\pi\) in chronological order (the spread sufficiently small as below and as in the radial estimate). These integrated fields are called screenings or root vertices.
\par\noindent\textbullet\quad If \(\epsilon=1\), the scaled rotated compensator at \(s_*\).
\par\noindent\textbullet\quad At 0 the single scaled vertex \(W_B^{\epsilon p}\), \(B=N h+\epsilon p\).

For each list of charges \(b_j,\sigma_j\) in this expression let \(d=\sum_j b_j\); \(d\in\mathcal L+\mathbb Z h\) and \(\sum_j\sigma_j=0\). Drop the term if \((p,d)\ne0\). Otherwise sum over
\[
 m=-K d+t p,\qquad n=A m=m-d,\qquad t\in\mathbb Z .
\]
Set \(n_{<j}=n+\sum_{l<j}b_l\), \(z_j=z(s_j)\). With the measures, compensation coefficients and sums just specified, the integrand (including angular partition functions) is
\begin{equation}
\begin{split}
 &\frac{Z_A}{Z_C}\, Q_L^{n^2/2}
     \prod_j\left[(2\pi/L)^{b_j^2/2} z_j^{(b_j,n_{<j})+b_j^2/2}\right]\\
 &\ \times\exp\!\left(2\pi i(\mu,n)+\sum_{j<k}[\pi i(b_k,K b_j)+2\pi i(b_j,K\sigma_k)]\right)\\
 &\ \times\prod_j\prod_{l\ge1}(1-Q_L^l)^{(A^l b_j,b_j)}
  \prod_{j<k}
  \left\{\prod_{l\ge0}(1-Q_L^l z_j/z_k)^{(A^l b_j,b_k)}
         \prod_{l\ge1}(1-Q_L^l z_k/z_j)^{(A^{-l} b_j,b_k)}\right\}.            \end{split}
\label{c4:Ang-2}
\end{equation}
Powers of \(1-\cdot\) use disk logs (polynomial for same-contour direct interactions). Two baseline identities will also be used: this same expression with only the screening factors, no \(B\) or compensator, equals 1 when summed/integrated, both for \(\mu\) as given and for \(\mu+p\).

We justify these identities first near \(Q_L=0\). Conjugate all operators by \(g^{-1}\), replacing \(\widetilde T\) by \(W_h^0\), rotating the compensator if present and factorizing the conjugated twist as above. Cyclically move \(\widehat U\) to the rightmost (earliest) position. For the baseline use no external vertices in the original angular trace, optionally inserting \(e^{2\pi i(p,P)}\) by \(J\) (also commuting with the zero modes); only zero momentum contributes before conjugation, giving 1.

There is a finiteness precaution in formal trace algebra here since \(H\) is indefinite. Work first at specified cut \(H\) and field powers (before scaling and fusing); other inserted operators preserve \(H\). In expressions with twist cut at \(C\), \(m\) before the twist satisfies \((1-C)m=d'\) with \(d'\) a fixed shift plus combinations from \(\alpha_6,\alpha_7,\alpha_0,\alpha_1\) (all of \(g\)). Thus \(m=-K d'\), with \(K\)-image Gram on these four vectors having leading minors \(9/8,5/4,11/8,1/2\), positive definite. Only finitely many \(m\) can contribute at fixed cut \(H\), including during the conjugation. Internal transitions have finite support from any basis state; this justifies expansions and cyclic moves on this side. With \(\widehat U\) expanded at the new cut, \((1-A)m=d\) is equivalent to \(m=-K d-(\alpha_0,m)p\) and imposes \((p,d)=0\); since \(K d\in\mathcal L^*\) it gives exactly the lattice sum stated. There are constants \(c_0>0,C_0\) such that
\[
 n^2=m^2\ \ge\ c_0\big(\textstyle\sum_b k_b^2+t^2\big)-C_0
\]
on possible terms (\(C_0\) may depend on fixed external charges). Indeed project \(K\) times the seven columns of \eqref{c4:Ang-1} off \(p\) orthogonally. The resulting Gram is given on columns \(1,2,5\) by
\[
 \begin{pmatrix}2&-1/3&2/3\\-1/3&10/9&1/9\\2/3&1/9&10/9\end{pmatrix}>0,
\]
with \(3,4,6,7\) reducing for the Gram to column 2, zero, \((\mathrm{col}\ 1+\mathrm{col}\ 5)/2\), column 5 respectively (multiply using \(K=\sum j C^j/8\)). Thus on nonnegative counts this bounds squares outside column 4; column 4 has nonzero \(p\)-pairing and is controlled by \((p,d)=0\). Fixed external shifts give only linear errors. The invariant part \(p[t-(p,Kd)/p^2]\) then also controls \(t\). This proves finiteness for the expanded moves. For the phases one has exactly
\[
 (d,K n)-n^2-(\alpha_0,m)(\alpha_0,K m)=0
\]
by \(n=Cm-(\alpha_0,m)\alpha_1\), \(m^2=n^2\). Thus all state-dependent cocycle, \(J\) and Weyl phases cancel except \(D\); the same with \(C\) only omits the Weyl term.

For clarity, the oscillator trace in a copy of index \(r'>0\) with twist \(Q_L^{r'} M\) (\(M=C,A\)) adds direct contraction
\(- (b_j,b_k)(z_j/z_k)^{r'}/r'\) in the exponential for creator at \(j\), annihilator at \(k>j\); also wrapped contractions for all \(j,k\) including equality,
\(- (M^l b_j,b_k)(Q_L^l z_j/z_k)^{r'}/r'\), \(l\ge1\). This follows by commuting to normal order and using
\[
 \operatorname{Tr}_{\rm Sym}\big[\operatorname{Sym}(D')\exp(c^\dagger(F))\exp(c(G))\big]
 =\det(1-D')^{-1}\exp(G D'(1-D')^{-1}F)
\]
for multiplication and differentiation with vector \(F\), covector \(G\), and matrix \(D'\). Diagonalize and extend formally; in one variable this is a sum of monomial diagonals with \(\sum_m x^m\binom{m}{j}=x^j/(1-x)^{j+1}\). Multiplying oscillator copies gives the products in \eqref{c4:Ang-2}, or the same products with \(C\) for the original unfused trace. All these manipulations so far are coefficientwise, with a separated pure cluster \(W_h^0\) and \(W_v^p\) as appropriate in place of the single \(B\) on the screening side.

These formal expansions can be evaluated by convergent sums for \(Q_L\) small and at the angular positions indicated, integrating for the root zero modes (their powers are integral, scaling exactly changes \(du/L\) into the stated scaled field measure). Indeed the exponent of \(Q_L\) in modulus from momentum heat and the \(z_j\) powers is the cyclic average of the successive momentum squares divided by 2, with weights angular gaps over \(2\pi\). Off the short screening cluster those squares equal \(m^2=n^2\) up to linear errors; within it they are bounded below by \(-C(\sum k_b^2+t^2+1)\). Choose the spread sufficiently small to keep a positive quadratic lower bound. Absolute expansion bounds for the oscillator factors are \(\exp(O((1+\sum k_b)^2))\) near \(Q_L=0\), since powers of the twists grow at most polynomially in \(l\), angles for distinct factors are separated and same-factor direct powers polynomial. This estimate at separated external positions follows also for sums of absolute coefficients weighted by the moduli (use the logarithm series); possible constants from external-external factors at tiny nonzero \(x\) do not affect the screening quadratic coefficient. Thus sums and root integrations recover the coefficientwise identity.

Now hold a fixed sufficiently small \(Q_L\) and take the divided cluster limit. On the transformed side the direct pair powers among the pure cluster are all \(h^2\), giving, since \(1-z_j/z_k\sim(2\pi i/L)(s_j-s_k)\), the additional scaling exactly into \((2\pi/L)^{B^2/2}\), up to a fixed unit phase. All in-cluster cocycles similarly give only one unit scalar. Other factors combine by bilinearity as in \eqref{c4:Ang-2}. The absolute screening bounds remain sufficient after division near this limit (screening and external rays stay separated), justifying it.

On the \(C\) side the divided expression has a Laurent expansion in \(x\) with at worst finite pole since in-cluster direct exponents in each summand differ from \(h^2\) by integers; all coefficients are real-analytic in \(L>0\) with \(u_*/L\) fixed. The negative Laurent coefficients vanish by the small-\(Q_L\) result. The sum of \eqref{c4:Ang-2} after transformation in the next calculation is also real-analytic on that interval, as will be proved there. This continues the equalities (including baseline cases) to real \(L>0\), using the radially summed expression for \eqref{c4:Ang-2}.

\subsection{Recovering the finite color sums}

The transformed expression now has a precise coefficientwise meaning and a continuation criterion. We identify its original finite-model observable before using the radial estimates to justify that continuation.

We identify the absolute confluent traces in the successive limits \(L\to\infty\) and (in the marked case) \(u_*\to+\infty\) from the \(C\) side. Here is the detailed pair limit. Initially separate the cluster on the indicated rays. For \(j<k\), \(a_l=(C^l b_j,b_k)\) has period 8 with zero mean, and put \(c'=i(s_j-s_k)/(2\pi)\). The oscillator pair products tend to
\begin{equation}
 \prod_{l=0}^7(2\sin(\pi(l+c')/8))^{a_l}
 = e^{-\pi i(b_k,K b_j)}
     \prod_{l=0}^7(1-q^l e^{-(s_j-s_k)/8})^{a_l}.                        \label{c4:Ang-3}
\end{equation}
In fact group offsets \(8m+l+c'\) from the first product and \(8m+8-l-c'\) from the second, \(m\ge0\). The weighted log-sums have cancelling constant and linear Taylor terms per group, with bounds \(O((m+1)^{-2})\) for \(m\ge1\) uniformly as \(Q_L\to1^-\) by differentiating \(\log(1-Q_L^y)\) twice. The first group is also locally uniform after extracting any \(c'=0\) singular power. Passing to the ordinary log-products gives the first formula by \(\sin(\pi z)/(\pi z)=\prod_{n\ge1}(1-z^2/n^2)\), constants canceling by zero mean. For instance this product follows by logarithmic derivative from \(\pi\cot(\pi z)=1/z+\sum_{n\ge1}2z/(z^2-n^2)\) (integrate \(\pi\cot(\pi t)/(t^2-z^2)\) on growing squares with real edges at half-integer levels); both sides of the product equal 1 at zero. Branches taken first at real \(0<c'<1\) and then with \(\Im c'>0\) give the second expression in \eqref{c4:Ang-3} with disk logs. Its phase cancels the conventional pair cocycle. Similarly the self product including vertex scaling tends to
\[
 d_C(b)=8^{-b^2/2}\prod_{l=1}^7 |1-q^l|^{(C^l b,b)/2};
\]
this follows by taking equal charges in the same formula, \(c'\to0^+\), dividing by the direct singular factor \((1-Q_L^{c'})^{b^2}\) while multiplying by \((2\pi/L)^{b^2}\), and taking the positive square root. The limits commute with Laurent coefficient extraction at \(x=0\) by the local uniform bounds just given after factoring the direct powers. Subsequently the separated compensator's pair interactions in the disk products tend to 1 analytically as \(u_*\to+\infty\). The extra \(\sigma\) pair phase in the marked case is \(e^{-2\pi i(v,K p)}\); the zero-mode powers and heat tend to 1 since \(m=-K\sum b_j,\ n=Cm\) are fixed per term.

From each cluster pair extract the disk product of \eqref{c4:Ang-3} with both charges \(h\); divided by the corresponding \((s_j-s_k)^{h^2}\) it tends to \(d_C(h)^2\). The remaining pair functions for \((b_j,b_k)=(h_i,h_{i'})\), \(1\le i,i'\le3\), are exactly the color pair functions \(g_{ii'}(r')\) with \(r'=\exp((s_k-s_j)/8)\). Indeed subtract \((C^l h,h)\) from the exponents: \(g_{11}=g_{33}=1,\ g_{22}=V_q(r'),\ g_{12}=g_{23}=V_q(r'/q),\ g_{13}=V_q(r'/q)V_q(r'/q^2),\ g_{i'i}(r')=g_{ii'}(1/r')\), \(V_q(t)=(1+t)^2/[(1-qt)(1-q^{-1}t)]\). For \((b_j,b_k)=(v,h_i)\) the resulting exponents of \(1-q^l r'\) are
\[
 (0,0,-1,-1,-1,1,1,1),\quad (0,0,0,0,0,0,0,0),\quad (0,1,1,1,-1,-1,-1,0)
\]
respectively, just the one-mark factors. Self weights combined with coefficients of \(\widetilde T\) at color sites are \(d_C(h)/w_-\) times \((w_-,w_0,w_+)\), since \(d_C(h_2)/d_C(h)=2/|1-q|\) and \(d_C(h_3)=d_C(h)\). Although individual residual pair terms may have poles, the summed color expressions here have the analytic limit (also by vanishing of the negative Laurent coefficients). Consequently the successive absolute-value limits of the confluent angular expressions are
\begin{equation}
 \begin{array}{ll}
 \epsilon=0:\quad d_C(h)^{N^2}|w_-|^{-N}|\tau_N|,\\
 \epsilon=1:\quad d_C(v)^2 d_C(h)^{N^2-1}|w_-|^{-(N-1)}|\tau^v_{N-1}|.
 \end{array}                                                              \label{c4:Ang-4}
\end{equation}

